% Options for packages loaded elsewhere
\PassOptionsToPackage{unicode}{hyperref}
\PassOptionsToPackage{hyphens}{url}
\documentclass[
  a4paper,
]{article}
\usepackage{xcolor}
\usepackage[margin=25mm]{geometry}
\usepackage{amsmath,amssymb}
\setcounter{secnumdepth}{-\maxdimen} % remove section numbering
\usepackage{iftex}
\ifPDFTeX
  \usepackage[T1]{fontenc}
  \usepackage[utf8]{inputenc}
  \usepackage{textcomp} % provide euro and other symbols
\else % if luatex or xetex
  \usepackage{unicode-math} % this also loads fontspec
  \defaultfontfeatures{Scale=MatchLowercase}
  \defaultfontfeatures[\rmfamily]{Ligatures=TeX,Scale=1}
\fi
\usepackage{lmodern}
\ifPDFTeX\else
  % xetex/luatex font selection
\fi
% Use upquote if available, for straight quotes in verbatim environments
\IfFileExists{upquote.sty}{\usepackage{upquote}}{}
\IfFileExists{microtype.sty}{% use microtype if available
  \usepackage[]{microtype}
  \UseMicrotypeSet[protrusion]{basicmath} % disable protrusion for tt fonts
}{}
\makeatletter
\@ifundefined{KOMAClassName}{% if non-KOMA class
  \IfFileExists{parskip.sty}{%
    \usepackage{parskip}
  }{% else
    \setlength{\parindent}{0pt}
    \setlength{\parskip}{6pt plus 2pt minus 1pt}}
}{% if KOMA class
  \KOMAoptions{parskip=half}}
\makeatother
\usepackage{longtable,booktabs,array}
\newcounter{none} % for unnumbered tables
\usepackage{calc} % for calculating minipage widths
% Correct order of tables after \paragraph or \subparagraph
\usepackage{etoolbox}
\makeatletter
\patchcmd\longtable{\par}{\if@noskipsec\mbox{}\fi\par}{}{}
\makeatother
% Allow footnotes in longtable head/foot
\IfFileExists{footnotehyper.sty}{\usepackage{footnotehyper}}{\usepackage{footnote}}
\makesavenoteenv{longtable}
\setlength{\emergencystretch}{3em} % prevent overfull lines
\providecommand{\tightlist}{%
  \setlength{\itemsep}{0pt}\setlength{\parskip}{0pt}}
\usepackage{bookmark}
\IfFileExists{xurl.sty}{\usepackage{xurl}}{} % add URL line breaks if available
\urlstyle{same}
\hypersetup{
  hidelinks,
  pdfcreator={LaTeX via pandoc}}

\author{}
\date{}

\begin{document}

\section{The Geometric Identity of the Zero-Square-Sum Constraint under
Scale
Invariance}\label{the-geometric-identity-of-the-zero-square-sum-constraint-under-scale-invariance}

\subsection{Isotropic cone, projective quadric, and intrinsic quantum
structure --- an expository
note}\label{isotropic-cone-projective-quadric-and-intrinsic-quantum-structure-an-expository-note}

\textbf{Author:} Noriaki Kihara \textbf{ORCID:} 0009-0004-6753-4020
\textbf{Date:} July 22, 2026 \textbf{Version DOI:}
10.5281/zenodo.21495306 \textbf{Concept DOI:} 10.5281/zenodo.21495305
\textbf{Position:} Expository note. It contains no new theorems; its
contribution is the synthesis of known mathematical results and their
placement onto one object.

\begin{center}\rule{0.5\linewidth}{0.5pt}\end{center}

\subsection{Abstract}\label{abstract}

Consider a system defined by exactly two postulates.

\begin{enumerate}
\def\labelenumi{\arabic{enumi}.}
\tightlist
\item
  States are complex vectors
  \(x=(x_1,\dots,x_N)\in\mathbb C^N\setminus\{0\}\) satisfying the
  zero-square-sum constraint \(\sum_{n=1}^N x_n^2=0\).
\item
  Only ratios are physical: \(x\) and \(\lambda x\)
  (\(\lambda\in\mathbb C^\ast\)) represent the same state.
\end{enumerate}

The purpose of this note is to identify the mathematical identity of the
state space these two postulates define, by citation of known results
alone. The conclusions are organized in five points.

\begin{enumerate}
\def\labelenumi{\arabic{enumi}.}
\tightlist
\item
  The solution set of postulate 1 is the complex isotropic (null) cone,
  and after projectivization by postulate 2 the state space is a quadric
  hypersurface in complex projective space --- the \textbf{projective
  quadric} \(Q_{N-2}\subset\mathbb{CP}^{N-1}\). This is a compact Kähler
  manifold {[}4{]}
\item
  By the general theory of geometric quantization, a system with compact
  phase space has, with no room for choice, \textbf{finite-dimensional
  state spaces and discrete spectra} {[}1,2,3{]}. Discreteness
  (quantumness) is not an additional hypothesis but a structure that
  follows automatically from the two postulates. The system is not
  merely quantizable --- it is born quantized
\item
  Writing \(x_n=q_n+ip_n\), the real part of the constraint is
  \(\sum q_n^2=\sum p_n^2\) (equipartition) and the imaginary part is
  \(\sum q_np_n=0\) --- the condition that the generator of scale
  transformations (the dilatation) of classical mechanics vanishes
  identically. Scale invariance (postulate 2) may not be an independent
  requirement but already present in the imaginary part of postulate 1
\item
  The space of polynomial functions on the cone is exactly the space of
  \textbf{harmonic polynomials} {[}6{]}, and harmonic polynomials give
  the spherical harmonics --- integer-labelled irreducible multiplets of
  the rotation group. The structure of ``harmonics'' appears not as a
  metaphor but as a theorem of the state space
\item
  Line bundles on compact projective manifolds are classified by
  integers (Chern classes) {[}4,5{]}, and the exact integrality of
  winding-type quantities is the integrality of cohomology itself. In
  particular, for \(N=3\) the projectivized isotropic cone is a conic
  isomorphic to the Riemann sphere \(\mathbb{CP}^1\), and null vectors
  can be written as squares of spinors (Cartan's construction {[}7{]}).
  The state space of the minimal nontrivial system is thus isomorphic to
  that of a qubit (the projective state space of spin 1/2)
\end{enumerate}

\begin{center}\rule{0.5\linewidth}{0.5pt}\end{center}

\subsection{1. Postulates and Notation}\label{postulates-and-notation}

We restate the two postulates.

\begin{quote}
\textbf{Postulate 1 (zero-square-sum constraint)}
\(x\in\mathbb C^N\setminus\{0\}\),
\(\displaystyle\sum_{n=1}^N x_n^2=0\).

\textbf{Postulate 2 (scale invariance)} Only ratios are physical:
\(x\sim\lambda x\) (\(\lambda\in\mathbb C^\ast\)).
\end{quote}

The solution set of postulate 1,

\[
\mathcal C=\Bigl\{x\in\mathbb C^N\setminus\{0\}\ :\ \sum_{n=1}^N x_n^2=0\Bigr\},
\]

is the \textbf{isotropic cone} (null cone) of the complex quadratic form
\(\sum x_n^2\). Over the reals there is no nontrivial solution (a
vanishing sum of real squares forces every component to vanish). A
nontrivial solution requires imaginary parts; in this sense the complex
structure is built into postulate 1.

For \(N=2\) the cone degenerates into the two lines \(x_2=\pm ix_1\),
and the projectivization is two points. Throughout we consider
\(N\ge3\), where smooth geometry appears.

The identification \(x\sim\lambda x\) of postulate 2 is the quotient by
the group \(\mathbb C^\ast\). Corresponding to the decomposition
\(\mathbb C^\ast=\{|\lambda|\}\times\{e^{i\varphi}\}\), the
unreadability of absolute scale (\(|\lambda|\)) and the unreadability of
common phase (\(e^{i\varphi}\)) are the modulus part and the argument
part of one and the same group.

\begin{center}\rule{0.5\linewidth}{0.5pt}\end{center}

\subsection{2. The State Space Is a Projective
Quadric}\label{the-state-space-is-a-projective-quadric}

The cone \(\mathcal C\) is invariant under \(\lambda x\)
(\(\lambda\in\mathbb C^\ast\)), so the quotient

\[
Q_{N-2}\ :=\ \mathcal C/\mathbb C^\ast\ =\ \bigl\{[x]\in\mathbb{CP}^{N-1}\ :\ \textstyle\sum_n x_n^2=0\bigr\}
\]

is well defined. It is a nondegenerate quadric hypersurface in complex
projective space \(\mathbb{CP}^{N-1}\) --- the \textbf{projective
quadric} --- a smooth compact complex manifold of complex dimension
\(N-2\) {[}4{]}.

\(\mathbb{CP}^{N-1}\) is a Kähler manifold with the Fubini--Study
metric, and the quadric, as a complex submanifold, inherits the Kähler
structure {[}4{]}. Therefore:

\begin{quote}
The state space defined by the two postulates is a compact Kähler
manifold.
\end{quote}

This is the identification from which the note departs. Everything below
is a known consequence of this one line.

\begin{center}\rule{0.5\linewidth}{0.5pt}\end{center}

\subsection{3. Compactness Supplies
Quantumness}\label{compactness-supplies-quantumness}

The general theory of geometric quantization (Kostant {[}1{]}, Souriau
{[}2{]}) teaches the following. A symplectic manifold \((M,\omega)\)
admits a prequantum line bundle \(L\) when \([\omega/2\pi]\) is an
integral cohomology class (the prequantization condition); on a Kähler
manifold, choosing the holomorphic polarization, the quantum state space
is constructed as the space of holomorphic sections
\(H^0(M,L^{\otimes k})\). If \(M\) is \textbf{compact}, this space is
\textbf{finite-dimensional} for every \(k\).

On the quadric the restriction of the Fubini--Study form represents the
hyperplane class (an integral class), so the prequantization condition
can be satisfied automatically. Therefore:

\begin{quote}
The system of the two postulates has, with no room for choice in the
quantization procedure, finite-dimensional state spaces and discrete
spectra.
\end{quote}

The classical explicit model showing that quantization on a compact
phase space forces discreteness and finite-dimensionality is Berezin's
example of the sphere {[}3{]} (quantization of the sphere = spin systems
= finite dimensions). In ordinary quantization, discreteness is obtained
by introducing \(\hbar\) from outside; in this system the compactness of
the (projectivized) cone supplies the discreteness. In this sense the
system is not merely quantizable --- it is \textbf{born quantized}.

\begin{center}\rule{0.5\linewidth}{0.5pt}\end{center}

\subsection{4. Real and Imaginary Parts of the Constraint: Equipartition
and
Dilatation}\label{real-and-imaginary-parts-of-the-constraint-equipartition-and-dilatation}

Writing \(x_n=q_n+ip_n\) (\(q_n,p_n\in\mathbb R\)),

\[
\sum_n x_n^2=\Bigl(\sum_n q_n^2-\sum_n p_n^2\Bigr)+2i\sum_n q_np_n,
\]

so postulate 1 decomposes into two real conditions:

\[
\text{real part:}\quad \sum_n q_n^2=\sum_n p_n^2,
\qquad
\text{imaginary part:}\quad \sum_n q_np_n=0.
\]

The real part, reading \(q\) as configuration and \(p\) as momentum, is
the harmonic-oscillator-type \textbf{equipartition} condition. The
imaginary part, \(D=\sum_n q_np_n\), is the classical \textbf{generator
of scale transformations (the dilatation)}. The imaginary part is thus
the condition that the generator of scale transformations vanishes
identically: the zero-square-sum constraint has, built in from the
start, a structure that kills the scale degree of freedom identically.

From this one observation follows: \textbf{scale invariance (postulate
2) may not be an independent requirement but already present in the
imaginary part of postulate 1.} This note records the observation and
does not decide the logical independence or dependence of the two
postulates.

\begin{center}\rule{0.5\linewidth}{0.5pt}\end{center}

\subsection{5. The Function Space on the Cone Is Harmonic
Polynomials}\label{the-function-space-on-the-cone-is-harmonic-polynomials}

Restricting polynomials on \(\mathbb C^N\) to the isotropic cone
\(\mathcal C\) is the same as quotienting by the ideal \((\sum x_n^2)\).
By the classical direct-sum decomposition

\[
\mathcal P_k=\mathcal H_k\oplus r^2\,\mathcal P_{k-2}
\qquad(\mathcal P_k:\text{homogeneous polynomials of degree }k,\ \mathcal H_k:\text{harmonic polynomials},\ r^2=\textstyle\sum x_n^2),
\]

the space of degree-\(k\) functions on the cone is exactly identified
with the harmonic polynomials \(\mathcal H_k\) (a consequence of the
vanishing of the symbol of the Laplacian on the cone; the standard
reference is Stein--Weiss {[}6{]}, Chapter IV). Its dimension is

\[
\dim\mathcal H_k=\binom{N+k-1}{k}-\binom{N+k-3}{k-2},
\]

and on the sphere \(\mathcal H_k\) gives the spherical harmonics --- the
irreducible multiplets of the rotation group with integer label \(k\).

The natural function system of this state space is therefore literally
\textbf{harmonics}. On this geometry, the word ``harmonic'' is not a
metaphor but a theorem of the function space.

\begin{center}\rule{0.5\linewidth}{0.5pt}\end{center}

\subsection{6. Integrality: Chern Classes and
Winding}\label{integrality-chern-classes-and-winding}

Complex line bundles on a compact projective manifold are classified by
the first Chern class \(c_1\in H^2(M;\mathbb Z)\) --- an
\textbf{integer-coefficient} cohomology class (Chern {[}5{]}; textbook
treatment in {[}4{]}). On the quadric,
\(H^2(Q_{N-2};\mathbb Z)\cong\mathbb Z\) (with the single exception
\(N=4\), where \(Q_2\cong\mathbb{CP}^1\times\mathbb{CP}^1\) and
\(H^2\cong\mathbb Z^2\)) {[}4{]}.

Consequently, the exact integrality of ``winding''-type quantities on
this state space (the number of turns of the phase along a closed loop,
the label of a line bundle) is not an approximate or statistical
property but the integrality of the cohomology of a compact projective
manifold itself. The integrality does not depend on any detail of
dynamics and is invariant under continuous deformations and scale
transformations.

\begin{center}\rule{0.5\linewidth}{0.5pt}\end{center}

\subsection{\texorpdfstring{7. The Minimal Nontrivial System \(N=3\):
Conic, Spinor,
Qubit}{7. The Minimal Nontrivial System N=3: Conic, Spinor, Qubit}}\label{the-minimal-nontrivial-system-n3-conic-spinor-qubit}

For \(N=3\), the projectivized isotropic cone is a smooth conic (quadric
curve) in \(\mathbb{CP}^2\), and a conic is isomorphic, as a rational
normal curve, to \(\mathbb{CP}^1\) --- the Riemann sphere {[}4{]}.

Moreover, null vectors can be written explicitly as squares of spinors
(Cartan's construction {[}7{]}). For a spinor \((a,b)\in\mathbb C^2\),
set

\[
x_1=a^2-b^2,\qquad x_2=i\,(a^2+b^2),\qquad x_3=-2ab;
\]

then \(x_1^2+x_2^2+x_3^2=0\) holds identically, and conversely every
isotropic vector of \(\mathbb C^3\) can be written in this form. The
ratio of the spinor components, \([a:b]\in\mathbb{CP}^1\), determines
the state.

The state space of the minimal nontrivial system is therefore the
Riemann sphere --- \textbf{isomorphic to the projective state space of a
qubit (a two-level system), i.e., of spin 1/2}. \(SU(2)\) acts on the
spinors, \(SO(3)\) acts on the isotropic vectors, and the double cover
\(SU(2)\to SO(3)\) is implemented by Cartan's construction itself. The
structures of spin and \(SU(2)\) are in a position to emerge from this
geometry without additional assumptions.

We note that the program of founding physical theory on projectivized
null structures and spinorization is the adjacent field that Penrose's
twistor theory {[}8{]} has developed for half a century. The components
of this note --- null cone, spinors, holomorphic structure --- are all
standard tools there.

\begin{center}\rule{0.5\linewidth}{0.5pt}\end{center}

\subsection{8. What This Note Does Not
Claim}\label{what-this-note-does-not-claim}

\begin{enumerate}
\def\labelenumi{\arabic{enumi}.}
\tightlist
\item
  All identifications in this note are \textbf{kinematic} (at the level
  of the state space). No correspondence with any particular dynamics,
  update rule, or physical system is asserted.
\item
  The statement of Section 4 that ``scale invariance is already present
  in the imaginary part of the constraint'' is an observation, not a
  decision on the logical relation of the two postulates.
\item
  The statement of Section 7 that ``the origin of spin is in a position
  to emerge'' is the indication of a mathematical isomorphism, not a
  claim of deriving physical spin.
\item
  This note contains no new theorems. Every mathematical claim belongs
  to the cited literature; the contribution is the placement and
  synthesis onto the system of the two postulates.
\end{enumerate}

\subsection{9. Classification of Claims}\label{classification-of-claims}

{\def\LTcaptype{none} % do not increment counter
\begin{longtable}[]{@{}
  >{\raggedright\arraybackslash}p{(\linewidth - 2\tabcolsep) * \real{0.5000}}
  >{\raggedright\arraybackslash}p{(\linewidth - 2\tabcolsep) * \real{0.5000}}@{}}
\toprule\noalign{}
\begin{minipage}[b]{\linewidth}\raggedright
Claim
\end{minipage} & \begin{minipage}[b]{\linewidth}\raggedright
Classification
\end{minipage} \\
\midrule\noalign{}
\endhead
\bottomrule\noalign{}
\endlastfoot
Projectivized isotropic cone = projective quadric, compact Kähler &
Application of known mathematics {[}4{]} \\
Quantization of compact phase space = finite-dimensional, discrete &
Application of known mathematics {[}1,2,3{]} \\
Real part = equipartition; imaginary part = vanishing dilatation &
Elementary computation \\
Scale invariance is present in the imaginary part of the constraint &
Observation (indication of a possibility) \\
Function space on the cone = harmonic polynomials = harmonics &
Application of known mathematics {[}6{]} \\
Exact integrality of winding = integrality of cohomology & Application
of known mathematics {[}4,5{]} \\
\(N=3\): conic \(\cong\mathbb{CP}^1\), squares of spinors, isomorphic to
a qubit & Application of known mathematics {[}4,7{]} \\
Physical origin of spin and \(SU(2)\) & Not claimed (indication of
position only) \\
\end{longtable}
}

\begin{center}\rule{0.5\linewidth}{0.5pt}\end{center}

\subsection{References}\label{references}

\begin{enumerate}
\def\labelenumi{\arabic{enumi}.}
\tightlist
\item
  Bertram Kostant, ``Quantization and unitary representations,'' in
  \emph{Lectures in Modern Analysis and Applications III}, Lecture Notes
  in Mathematics 170, 87--208, Springer, 1970. DOI:
  \href{https://doi.org/10.1007/BFb0079068}{10.1007/BFb0079068}.
\item
  Jean-Marie Souriau, \emph{Structure of Dynamical Systems: A Symplectic
  View of Physics}, Progress in Mathematics 149, Birkhäuser, 1997
  (French original: \emph{Structure des systèmes dynamiques}, Dunod,
  1970). DOI:
  \href{https://doi.org/10.1007/978-1-4612-0281-3}{10.1007/978-1-4612-0281-3}.
\item
  Felix A. Berezin, ``General concept of quantization,''
  \emph{Communications in Mathematical Physics}, 40, 153--174, 1975.
  DOI: \href{https://doi.org/10.1007/BF01609397}{10.1007/BF01609397}.
\item
  Phillip Griffiths and Joseph Harris, \emph{Principles of Algebraic
  Geometry}, Wiley, 1978 (Wiley Classics Library, 1994). DOI:
  \href{https://doi.org/10.1002/9781118032527}{10.1002/9781118032527}.
\item
  Shiing-Shen Chern, ``Characteristic classes of Hermitian manifolds,''
  \emph{Annals of Mathematics}, 47(1), 85--121, 1946. DOI:
  \href{https://doi.org/10.2307/1969037}{10.2307/1969037}.
\item
  Elias M. Stein and Guido Weiss, \emph{Introduction to Fourier Analysis
  on Euclidean Spaces}, Princeton University Press, 1971, Chapter IV.
  ISBN: 978-0-691-08078-9.
\item
  Élie Cartan, \emph{The Theory of Spinors}, Hermann, 1966 (French
  original: \emph{Leçons sur la théorie des spineurs}, 1938; Dover
  reprint 1981). ISBN: 978-0-486-64070-9.
\item
  Roger Penrose, ``Twistor algebra,'' \emph{Journal of Mathematical
  Physics}, 8(2), 345--366, 1967. DOI:
  \href{https://doi.org/10.1063/1.1705200}{10.1063/1.1705200}.
\end{enumerate}

\end{document}
